\documentclass[11pt]{article}

% ============================================================
% Packages
% ============================================================
\usepackage[utf8]{inputenc}
\usepackage[margin=1in]{geometry}
\usepackage{amsmath,amssymb,amsthm}
\usepackage{graphicx}
\usepackage{booktabs}
\usepackage{xcolor}
\usepackage{hyperref}
\usepackage{tikz}
\usepackage{enumitem}
\usepackage{pgfplots}

\pgfplotsset{compat=1.18}

% ============================================================
% Document Formatting
% ============================================================
\setlength{\parskip}{0.6em}
\setlength{\parindent}{0pt}

% ============================================================
% List Formatting
% ============================================================
\setlist[enumerate]{
    leftmargin=2em,
    itemsep=0.25em
}

% ============================================================
% Solution Environment
% ============================================================
\newenvironment{solution}
    {\par\noindent\textbf{Solution.}\par}
    {\par}

% ============================================================
% Link Styling
% ============================================================
\hypersetup{
    colorlinks=true,
    linkcolor=blue,
    urlcolor=cyan
}

% ============================================================
% Metadata
% ============================================================
\title{Ordinary Differential Equations}
\author{Chris Ritter}
\date{August 27, 2026}

% ============================================================
% Document
% ============================================================

\begin{document}

\maketitle
\tableofcontents
\newpage

% ============================================================
\section{Linear First-Order Equations}
% ============================================================

A first-order linear differential equation has the form

\[
y' + p(t)y = g(t).
\]

The goal is to find an integrating factor $\mu(t)$ such that
multiplying the entire equation by $\mu(t)$ produces the derivative
of a product.

% ------------------------------------------------------------
\subsection{Integrating Factor}
% ------------------------------------------------------------

We choose

\[
\mu(t)
=
e^{\int p(t)\,dt}.
\]

Multiplying

\[
y' + p(t)y = g(t)
\]

by $\mu(t)$ gives

\[
\mu y' + p(t)\mu y = \mu g.
\]

We want the left-hand side to be

\[
\frac{d}{dt}(\mu y).
\]

By the product rule,

\[
\frac{d}{dt}(\mu y)
=
\mu' y + \mu y'.
\]

Therefore, we need

\[
\mu' = p(t)\mu.
\]

This is satisfied by

\[
\mu(t)=e^{\int p(t)\,dt}.
\]

Thus,

\[
\frac{d}{dt}(\mu y)=\mu g.
\]

Integrating both sides gives

\[
\mu y
=
\int \mu(t)g(t)\,dt+C.
\]

Therefore,

\[
\boxed{
y(t)
=
\frac{1}{\mu(t)}
\left(
\int \mu(t)g(t)\,dt+C
\right)
}.
\]

% ------------------------------------------------------------
\subsection{Example: Linear Equation}
% ------------------------------------------------------------

\textbf{Solve}

\[
\frac{dy}{dt}+\frac12y=t.
\]

This is in the form

\[
y'+p(t)y=g(t),
\]

where

\[
p(t)=\frac12,
\qquad
g(t)=t.
\]

The integrating factor is

\[
\mu(t)
=
e^{\int \frac12\,dt}
=
e^{t/2}.
\]

Multiply the differential equation by $e^{t/2}$:

\[
e^{t/2}y'
+
\frac12e^{t/2}y
=
te^{t/2}.
\]

The left-hand side is the derivative of a product:

\[
\frac{d}{dt}
\left(
e^{t/2}y
\right)
=
te^{t/2}.
\]

Integrate:

\[
e^{t/2}y
=
\int te^{t/2}\,dt+C.
\]

Using integration by parts,

\[
\int te^{t/2}\,dt
=
e^{t/2}(2t-4).
\]

Therefore,

\[
e^{t/2}y
=
e^{t/2}(2t-4)+C.
\]

Divide by $e^{t/2}$:

\[
\boxed{
y(t)=2t-4+Ce^{-t/2}
}.
\]

% ============================================================
\section{Separable Equations}
% ============================================================

A first-order differential equation

\[
\frac{dy}{dx}=f(x,y)
\]

is called \textbf{separable} if it can be written in the form

\[
M(x)\,dx+N(y)\,dy=0.
\]

Equivalently, we can write

\[
N(y)\frac{dy}{dx}=-M(x).
\]

The important feature is that the $x$-dependent terms and
$y$-dependent terms can be separated.

% ------------------------------------------------------------
\subsection{General Method}
% ------------------------------------------------------------

Suppose we have

\[
M(x)\,dx+N(y)\,dy=0.
\]

Let $H_1$ and $H_2$ be antiderivatives satisfying

\[
H_1'(x)=M(x)
\]

and

\[
H_2'(y)=N(y).
\]

Then

\[
\frac{d}{dx}H_1(x)
+
\frac{d}{dx}H_2(y)=0.
\]

By the chain rule,

\[
\frac{d}{dx}H_2(y)
=
H_2'(y)\frac{dy}{dx}
=
N(y)\frac{dy}{dx}.
\]

Thus,

\[
\frac{d}{dx}
\left(
H_1(x)+H_2(y)
\right)
=0.
\]

Therefore,

\[
\boxed{
H_1(x)+H_2(y)=C
}.
\]

% ------------------------------------------------------------
\subsection{Initial Value Problems}
% ------------------------------------------------------------

Suppose an initial condition

\[
y(x_0)=y_0
\]

is given.

Since

\[
H_1(x)+H_2(y)=C,
\]

substituting $(x_0,y_0)$ gives

\[
C
=
H_1(x_0)+H_2(y_0).
\]

Therefore, the solution to the IVP can be written as

\[
\boxed{
H_1(x)+H_2(y)
=
H_1(x_0)+H_2(y_0)
}.
\]

This gives the solution implicitly. If possible, we may then solve
for $y$ explicitly.

% ------------------------------------------------------------
\subsection{Example: Separable IVP}
% ------------------------------------------------------------

\textbf{Solve the IVP}

\[
\frac{dy}{dx}
=
\frac{3x^2+4x+2}{2(y-1)},
\qquad
y(0)=-1.
\]

First separate the variables:

\[
2(y-1)\,dy
=
(3x^2+4x+2)\,dx.
\]

Integrate both sides:

\[
\int 2(y-1)\,dy
=
\int(3x^2+4x+2)\,dx.
\]

The left-hand side becomes

\[
(y-1)^2,
\]

while the right-hand side becomes

\[
x^3+2x^2+2x+C.
\]

Thus,

\[
(y-1)^2
=
x^3+2x^2+2x+C.
\]

Use the initial condition

\[
y(0)=-1.
\]

Then

\[
(-1-1)^2=C,
\]

so

\[
C=4.
\]

Therefore,

\[
(y-1)^2
=
x^3+2x^2+2x+4.
\]

Taking the square root gives

\[
y-1
=
\pm
\sqrt{x^3+2x^2+2x+4}.
\]

Since $y(0)=-1$, we must choose the negative branch:

\[
\boxed{
y(x)
=
1-
\sqrt{x^3+2x^2+2x+4}
}.
\]

This is the solution satisfying the initial condition.

% ============================================================
\section{Homogeneous Equations}
% ============================================================

A first-order differential equation

\[
\frac{dy}{dx}=f(x,y)
\]

is called \textbf{homogeneous} if the right-hand side can be
written as a function of the ratio $y/x$:

\[
\boxed{
\frac{dy}{dx}=g\left(\frac{y}{x}\right)
}.
\]

The important observation is that a homogeneous equation can be
converted into a separable equation using the substitution

\[
\boxed{
v=\frac{y}{x}
}.
\]

Equivalently,

\[
y=vx.
\]

% ------------------------------------------------------------
\subsection{The Homogeneous Substitution}
% ------------------------------------------------------------

Starting with

\[
y=vx,
\]

differentiate both sides with respect to $x$:

\[
\frac{dy}{dx}
=
v+x\frac{dv}{dx}.
\]

Therefore, whenever

\[
\frac{dy}{dx}
=
g\left(\frac{y}{x}\right),
\]

we substitute

\[
\frac{y}{x}=v
\]

and

\[
\frac{dy}{dx}
=
v+x\frac{dv}{dx}.
\]

This gives

\[
v+x\frac{dv}{dx}=g(v).
\]

Rearranging,

\[
x\frac{dv}{dx}=g(v)-v.
\]

Hence,

\[
\frac{dv}{g(v)-v}
=
\frac{dx}{x}.
\]

The equation is now separable.

% ------------------------------------------------------------
\subsection{Example: Homogeneous Equation}
% ------------------------------------------------------------

\textbf{Solve}

\[
\frac{dy}{dx}
=
\frac{x^2+xy+y^2}{x^2}.
\]

First rewrite the right-hand side:

\[
\frac{x^2+xy+y^2}{x^2}
=
1+\frac{y}{x}
+\left(\frac{y}{x}\right)^2.
\]

Therefore,

\[
\frac{dy}{dx}
=
1+\frac{y}{x}
+\left(\frac{y}{x}\right)^2.
\]

This is homogeneous because the right-hand side depends only on
$y/x$.

Use the substitution

\[
v=\frac{y}{x},
\]

so that

\[
y=vx.
\]

Differentiate:

\[
\frac{dy}{dx}
=
v+x\frac{dv}{dx}.
\]

Substitute into the differential equation:

\[
v+x\frac{dv}{dx}
=
1+v+v^2.
\]

Cancel $v$ from both sides:

\[
x\frac{dv}{dx}
=
1+v^2.
\]

Separate the variables:

\[
\frac{dv}{1+v^2}
=
\frac{dx}{x}.
\]

Integrate:

\[
\int\frac{dv}{1+v^2}
=
\int\frac{dx}{x}.
\]

Therefore,

\[
\arctan(v)
=
\ln|x|+C.
\]

Substitute back

\[
v=\frac{y}{x}:
\]

\[
\arctan\left(\frac{y}{x}\right)
=
\ln|x|+C.
\]

This is an implicit form of the solution.

Solving for $y/x$,

\[
\frac{y}{x}
=
\tan(\ln|x|+C).
\]

Thus,

\[
\boxed{
y
=
x\tan(\ln|x|+C)
}.
\]

% ============================================================
\section{Summary of Methods}
% ============================================================

\subsection{Linear First-Order Equations}

For

\[
y'+p(t)y=g(t),
\]

use the integrating factor

\[
\boxed{
\mu(t)=e^{\int p(t)\,dt}
}.
\]

Then

\[
\boxed{
\frac{d}{dt}(\mu y)=\mu g
}
\]

and therefore

\[
\boxed{
y
=
\frac{1}{\mu}
\left(
\int\mu g\,dt+C
\right)
}.
\]

% ------------------------------------------------------------
\subsection{Separable Equations}

If an equation can be written as

\[
M(x)\,dx+N(y)\,dy=0,
\]

then integrate:

\[
\boxed{
\int M(x)\,dx+\int N(y)\,dy=C
}.
\]

For an IVP

\[
y(x_0)=y_0,
\]

use the initial condition to determine $C$.

% ------------------------------------------------------------
\subsection{Homogeneous Equations}

If

\[
\frac{dy}{dx}
=
g\left(\frac{y}{x}\right),
\]

use

\[
\boxed{
v=\frac{y}{x}
}
\]

so that

\[
\boxed{
y=vx
}
\]

and

\[
\boxed{
\frac{dy}{dx}
=
v+x\frac{dv}{dx}
}.
\]

The resulting equation should be separable.

% ============================================================
\end{document}
